% Part: first-order-logic
% Chapter: syntax-and-semantics
% Section: unique-readability

\documentclass[../../../include/open-logic-section]{subfiles}

\begin{document}

\olfileid{fol}{syn}{unq}

\olsection{يوازينے لوست}

\begin{explain}
موږ چې !!{formula}s په کومه طريقه تعريف کړي دي، هغه دا تضمينوي چې هر
!!{formula} يو \emph{يوازينے لوست} لري؛ يعنې د !!{formula}s د جوړولو
زموږ د قاعدو له مخې ئې د جوړولو په اصل کښې يوازې يوه طريقه شته، او د
``تفسيرولو'' ئې هم يوازې يوه طريقه شته۔ که داسې نۀ واے، مبهم
!!{formula}s به مو لرل؛ يعنې داسې !!{formula}s چې له يوه څخه زيات
لوستونه يا تفسيرونه لري---او څرګنده ده چې موږ ترې ډډه کول غواړو۔ خو
تر دې مهمه دا ده چې له دې خاصيت پرته به زموږ ډېری راتلونکي تعريفونه
او ثبوتونه کار ورنۀ کړي۔
\textbf{د ژباړې د نسخې سمون (OLFOL-014):} سرچينه په همدې جمله کښې
د «تفسير» د انګريزي کلمې املا ناسمه ليکي؛ ژباړه ئې سمه معنا راوړي۔

د دې خبرې د روښانولو ښايي غوره لار دا وي چې وګورو، که د !!{formula}s
د جوړولو داسې ناسمې قاعدې مو ورکړې واے چې يوازينے لوست نۀ تضمينوي،
څه به پېښ شوي وو۔ مثلاً کېدے شي د نښلوونکو د جوړولو په قاعدو کښې مو
قوسونه هېر کړي او دا مو روا ګڼلي واے:
\begin{quote}
که $!A$ او $!B$ !!{formula}s وي، نو $!A \lif !B$ هم فارمول دے۔
\end{quote}
له يوه اټومي فارمول $!D$ څخه په پيلولو به دې قاعدې د $!D \lif !D$
جوړول را ته روا کړي وو۔ له دې او $!D$ څخه به مو $!D \lif !D \lif
!D$ ترلاسه کړے و۔ خو دا په دوو طريقو جوړېدے شي:
\begin{enumerate}
\item $!D$ د $!A$ په توګه، او $!D \lif !D$ د $!B$ په توګه اخلو۔
\item $!A$ د $!D \lif !D$ په توګه اخلو، او $!B$ پخپله~$!D$ دے۔
\end{enumerate}
په همدې ډول د !!{formula}~$!D \lif !D \lif !D$ د ``لوستلو'' دوه
طريقې شته۔ دا د $!B \lif !C$ بڼه لري چې $!B$ پکښې $!D$ او $!C$
پکښې $!D \lif !D$ دے، خو \emph{همدارنګه} د $!B \lif !C$ بڼه لري
چې $!B$ پکښې $!D \lif !D$ او $!C$ پکښې~$!D$ دے۔

که داسې وشي، زموږ تعريفونه به تل کار نۀ کوي۔ مثلاً کله چې د يوه فارمول
!!{main operator} تعريفوو، وايو: د $!B \lif !C$ بڼه لرونکي فارمول
!!{main operator} د~$\lif$ همدا ښودل شوے مورد دے۔ خو که فارمول $!D
\lif !D \lif !D$ له $!B \lif !C$ سره په پورته يادو شوو دوو بېلو
طريقو برابرولے شو، نو په يوه حالت کښې د $\lif$ لومړے مورد او په بل
حالت کښې ئې دويم مورد د !!{main operator} په توګه ترلاسه کوو۔ خو موږ
غواړو چې !!{main operator} د !!{formula} يوه \emph{تابع} وي؛ يعنې هر
!!{formula} بايد کټ مټ د !!{main
  operator} يو مورد ولري۔
\end{explain}

\begin{lem}
په !!a{formula}~$!A$ کښې د کيڼو او ښيو قوسونو شمېر سره برابر دے۔
\end{lem}

\begin{proof}
دا د $!A$ د جوړېدو پر طريقه په استقرا ثابتوو۔ دې ته دوه څيزونه پکار
دي: (الف) لومړے بايد ثابته کړو چې ټول اټومي فارمولونه ياد خاصيت لري
(د استقرا بنسټ)۔ (ب) بيا بايد ثابته کړو چې کله له راکړل شوو فارمولونو
څخه نوي فارمولونه جوړوو، که زاړه فارمولونه ياد خاصيت ولري، نو نوي ئې
هم لري۔

$l(!A)$ دې په~$!A$ کښې د کيڼو قوسونو شمېر وي او $r(!A)$ دې د ښيو
قوسونو شمېر وي؛ $l(t)$ او $r(t)$ دې په همدې ډول په يوه ترم~$t$ کښې
د کيڼو او ښيو قوسونو شمېرونه وي۔

\begin{prob}
    ثابته کړئ چې د هر ترم~$t$ دپاره $l(t) = r(t)$۔
\end{prob}

\begin{enumerate}
\tagitem{prvFalse}{\indcase{!A}{\lfalse}{$\indfrm$ $0$ کيڼ او $0$
    ښي قوسونه لري۔}}{}

\tagitem{prvTrue}{\indcase{!A}{\ltrue}{$\indfrm$ $0$ کيڼ او $0$
    ښي قوسونه لري۔}}{}

\item \indcase{!A}{\Atom{R}{t_1,\dots,t_n}}{$l(\indfrm) = 1 + l(t_1) +
  \dots + l(t_n) = 1 + r(t_1) + \dots + r(t_n) = r(\indfrm)$۔ دلته
  هغه حقيقت کاروو چې د تمرين په توګه پرېښودل شوے دے: د هر ترم~$t$
  دپاره $l(t) = r(t)$۔}

\item \indcase{!A}{\eq[t_1][t_2]}{$l(\indfrm) = l(t_1) + l(t_2) =
  r(t_1) + r(t_2) = r(\indfrm)$۔}

\tagitem{prvNot}{\indcase{!A}{\lnot !B}{د استقرايي فرض له مخې
    $l(!B) = r(!B)$۔ نو $l(\indfrm) = l(!B) = r(!B) =
    r(\indfrm)$۔}}{}

\item \indcase{!A}{(!B \ast !C)}{د استقرايي فرض له مخې $l(!B) =
  r(!B)$ او $l(!C) = r(!C)$۔ نو $l(\indfrm) = 1 + l(!B) + l(!C) =
  1 + r(!B) + r(!C) = r(\indfrm)$۔}

\tagitem{prvAll}{\indcase{!A}{\lforall[x][!B]}{د استقرايي فرض له مخې
    $l(!B) = r(!B)$۔ نو $l(\indfrm) = l(!B) = r(!B) =
    r(\indfrm)$۔}}{}

% Print case for \lexists if prvEx and not prvAll
\tagitem{prvAll,notprvEx}{}{\indcase{!A}{\lexists[x][!B]}{د استقرايي
    فرض له مخې $l(!B) = r(!B)$۔ نو $l(\indfrm) = l(!B) = r(!B) =
    r(\indfrm)$۔}}

% Just say similarly if prvEx and prvAll
\tagitem{notprvAll,notprvEx}{}{\indcase{!A}{\lexists[x][!B]}{په همدې ډول۔}}
\end{enumerate}
\end{proof}

\begin{defn}[حقيقي مخکښ]
د سمبولونو تار $!B$ د سمبولونو د تار~$!A$ \emph{حقيقي مخکښ} دے، که
له $!B$ سره د سمبولونو د يوه ناتش تار نښلول~$!A$ ترلاسه کړي۔
\end{defn}

\begin{lem}\ollabel{lem:no-prefix}
که $!A$ !!a{formula} وي او $!B$ د $!A$ حقيقي مخکښ وي، نو $!B$
!!a{formula} نۀ دے۔
\end{lem}

\begin{proof}
تمرين۔
\end{proof}

\begin{prob}
\olref[fol][syn][unq]{lem:no-prefix} ثابته کړئ۔
\end{prob}

\begin{prop}
\ollabel{prop:unique-atomic}
که $!A$ يو اټومي !!{formula} وي، نو له لاندې شرطونو څخه يو، او يوازې
يو، پوره کوي۔
\begin{enumerate}
\tagitem{prvFalse}{$!A \ident \lfalse$۔}{}
\tagitem{prvTrue}{$!A \ident \ltrue$۔}{}
\item $!A \ident \Atom{R}{t_1,\dots,t_n}$، چې $R$ يو $n$-ځايه
  !!{predicate}، $t_1$، \dots، $t_n$ ترمونه، او هر يو $R$،
  $t_1$، \dots، $t_n$ په يوازيني ډول ټاکل شوي دي۔
\item $!A \ident \eq[t_1][t_2]$، چې $t_1$ او $t_2$ په يوازيني ډول
  ټاکل شوي ترمونه دي۔
\end{enumerate}
\end{prop}

\begin{proof}
تمرين۔
\end{proof}

\begin{prob}
\olref[fol][syn][unq]{prop:unique-atomic} ثابته کړئ۔ (لارښوونه: د
ترمونو دپاره د \olref[fol][syn][unq]{lem:no-prefix} يوه بڼه بيان او
ثابته کړئ۔)
\end{prob}

\begin{prop}[يوازينے لوست]
هر !!{formula} له لاندې شرطونو څخه يو، او يوازې يو، پوره کوي۔
\begin{enumerate}
\item $!A$ اټومي دے۔

\tagitem{prvNot}{$!A$ د $\lnot !B$ بڼه لري۔}{}

\tagitem{prvAnd}{$!A$ د $(!B \land !C)$ بڼه لري۔}{}

\tagitem{prvOr}{$!A$ د $(!B \lor !C)$ بڼه لري۔}{}

\tagitem{prvIf}{$!A$ د $(!B \lif !C)$ بڼه لري۔}{}

\tagitem{prvIff}{$!A$ د $(!B \liff !C)$ بڼه لري۔}{}

\tagitem{prvAll}{$!A$ د $\lforall[x][!B]$ بڼه لري۔}{}

\tagitem{prvEx}{$!A$ د $\lexists[x][!B]$ بڼه لري۔}{}
\end{enumerate}
سربېره پر دې، په هر حالت کښې $!B$، يا $!B$ او $!C$، په يوازيني ډول
ټاکل شوي دي۔ يعنې مثلاً د $!B$، $!C$ او $!B'$، $!C'$ داسې بېلې جوړې
نشته چې $!A$ په يو وخت د دواړو بڼه ولري:
\iftag{prvIf}{$(!B \lif !C)$ او $(!B' \lif !C')$۔}{
    \iftag{prvOr}{$(!B \lor !C)$ او $(!B' \lor !C')$۔}{
      \iftag{prvAnd}{$(!B \land !C)$ او $(!B' \land !C')$۔}{}}}
\end{prop}

\begin{proof}
د جوړولو قاعدې غواړي چې که !!a{formula} اټومي نۀ وي، بايد په
پرانيستي قوس~(، \iftag{prvNot}{$\lnot$،}{} يا په يوه سور پيل شي۔ بل
خوا، هر !!{formula} چې له لاندې سمبولونو څخه په يوه پيلېږي، بايد
اټومي وي: !!a{predicate}، !!a{function}، !!a{constant}
\iftag{prvFalse}{، $\lfalse$}{}\iftag{prvTrue}{، $\ltrue$}{}۔

نو په اصل کښې يوازې دا ښودل پکار دي چې که $!A$ د $(!B \ast !C)$
او هم د $(!B' \mathbin{\ast'} !C')$ بڼه لري، نو $!B \ident !B'$،
$!C \ident !C'$، او $\ast = {\ast'}$۔

نو فرض کړئ چې هم $!A \ident (!B \ast !C)$ او هم $!A \ident (!B'
\mathbin{\ast'} !C')$۔ اوس يا $!B \ident !B'$ دے، يا نۀ دے۔ که وي،
نو څرګنده ده چې $\ast = {\ast'}$ او $!C \ident !C'$، ځکه بيا دا د
$!A$ هغه فرعي تارونه دي چې له يوه ځايه پيلېږي او يو شان اوږدوالے لري۔
بل حالت $!B \not\ident !B'$ دے۔ څرنګه چې $!B$ او $!B'$ دواړه د
$!A$ هغه فرعي تارونه دي چې له يوه ځايه پيلېږي، يو بايد د بل حقيقي
مخکښ وي۔ خو دا د \olref{lem:no-prefix} له مخې ناشونې ده۔
\end{proof}


\end{document}
